\documentclass[11pt,reqno]{amsart}

\usepackage[a4paper,margin=1.15in]{geometry}
\usepackage[T1]{fontenc}
\usepackage{lmodern}
\usepackage{amsmath,amssymb,mathtools}
\usepackage{microtype}
\usepackage{enumitem}
\usepackage[hidelinks]{hyperref}
\raggedbottom
\hypersetup{
  pdftitle={Finite Pauli Densities and Exact Compatibility},
  pdfauthor={}
}

\numberwithin{equation}{section}

\newtheorem{theorem}{Theorem}[section]
\newtheorem{proposition}[theorem]{Proposition}
\newtheorem{lemma}[theorem]{Lemma}
\theoremstyle{definition}
\newtheorem{definition}[theorem]{Definition}
\theoremstyle{remark}
\newtheorem{remark}[theorem]{Remark}

\newcommand{\cA}{\mathcal A}
\newcommand{\cH}{\mathcal H}
\newcommand{\cM}{\mathcal M}
\newcommand{\cG}{\mathcal G}
\newcommand{\cS}{\mathcal S}
\newcommand{\bbC}{\mathbb C}
\newcommand{\id}{\mathbf 1}
\newcommand{\Tr}{\operatorname{Tr}}
\newcommand{\rhoF}{\rho_F}
\newcommand{\PF}{P_F}

\title[Finite Pauli density realization]{Finite Pauli Densities and Exact Compatibility\\
\large A checked finite-volume realization for the stabilizer finite-fiber proof}
\author{}
\date{August 3, 2026}

\begin{document}

\begin{abstract}
This note records Investigation~I003, the finite-volume realization layer for
the finite-fiber formalization.  It machine-checks explicit qubit Pauli
matrices, tensor-product trace orthogonality, finite signed stabilizer group
averages, their trace-normalized density matrices, positivity and normalization
of the associated finite functional, and its traciality after compression to
the finite code space.  It also machine-checks finite partial trace for
signed Pauli sums.  Its generic partial-trace layer concerns normalized signed
Pauli expressions: an explicit survivor equality records one required part of
the conditional actual-density bridge described below, rather than granting
arbitrary term families the density properties.

These results establish the actual finite-group density calculation and the
algebraic partial-trace calculation in the current Overleaf proof.  They also
establish the conditional bridge from actual nested finite stabilizer groups to
the generic expression data: a supplied tensor-coordinate decomposition,
injective local-group inclusion, and exact survivor iff-condition imply
compatibility of the positive, trace-one local densities.  They do not derive
that contract from the countable-lattice local-group definition.  The checked
initial/final $\operatorname{Fin}$ coordinate split now yields typed unital
star-algebra maps between every ordered pair of finite volumes.  Their
identity, composition, injectivity, and operator-norm preservation are
checked, and they form a directed system.  At a target written as an
explicitly added block, the endpoint map is
also checked equal to the earlier length-indexed recursion, so its
signed-Pauli identity-padding calculation transfers without an implicit
identification.  The compatible values descend to its module direct limit as a
complex-linear functional.  This does not construct a noncommutative
star-algebra colimit on that quotient, the norm-closed quasi-local algebra or
a state on it, or remove the explicit GNS-orbit and corner-traciality
hypotheses in the central-fiber theorem of Investigation~I002.
\end{abstract}

\maketitle

\section{Contract and relation to the manuscript}
\label{sec:contract}

The source proof is the live manuscript
\path{writing/overleaf_prl_stabilizer_region_algebras/main.tex}.
Its finite local construction is in the appendix ``Infinite Stabilizer
Systems'': local groups and the code projector occur in lines 881--902, the
density formula in lines 903--914, and the partial-trace compatibility
calculation in lines 916--949.  The support-corner conclusion consuming the
GNS data is formalized separately in Investigation~I002.

The present I003 Lean source lies in the following package directory:
\begin{center}
  \path{formalization_finite_fibers/FiniteFiber/}
\end{center}
Its modules are:
\begin{itemize}[noitemsep,topsep=0.2em]
  \item \path{FinitePauli.lean};
  \item \path{FiniteQubitMatrix.lean};
  \item \path{FiniteDensity.lean};
  \item \path{FinitePauliStabilizer.lean};
  \item \path{FiniteLocalPauli.lean};
  \item \path{FiniteNestedPauli.lean};
  \item \path{FiniteLocalGroupSupport.lean};
  \item \path{ProjectivePauliLattice.lean};
  \item \path{FullPauliLattice.lean};
  \item \path{FiniteTensorMatrixChain.lean};
  \item \path{FiniteNestedTensorState.lean};
  \item \path{FiniteStabilizerStateChain.lean};
  \item \path{SequentialLocalFunctional.lean};
  \item \path{FinitePauliGNS.lean}.
\end{itemize}
No
theorem below uses a project-local axiom,
placeholder, or an opaque stabilizer-state construction.

\section{Finite signed-Pauli stabilizers}
\label{sec:finite_stabilizers}

For $n$ qubits, the Lean development represents the computational basis by a
finite type of cardinality $2^n$ and defines the four Hermitian one-qubit
matrices $id,X,Y,Z$.  It constructs the recursively parenthesized tensor
product $P(a)$ of a label function $a:\operatorname{Fin}(n)\to
\{\id,X,Y,Z\}$ and proves
\begin{equation}
  \Tr(P(a))=0
  \quad\text{if some }a(i)\neq\id,
  \qquad
  \Tr(P(\id,\ldots,\id))=2^n.
  \label{eq:pauli_trace_orthogonality}
\end{equation}

The finite stabilizer interface is not an unstructured assertion that a set
is a stabilizer group.  It contains a group homomorphism
\begin{equation}
  r:G\longrightarrow \operatorname{Mat}_{2^n}(\bbC),
  \label{eq:finite_representation}
\end{equation}
an explicit sign and Pauli label for every $g\in G$, injectivity of $r$, and
the condition $r(g)\neq-\id$ for every $g$.  Thus a nonidentity $g$ cannot be
the all-identity string with either sign, and the checked conclusion is
\begin{equation}
  g\neq e\quad\Longrightarrow\quad\Tr(r(g))=0.
  \label{eq:nonidentity_traceless}
\end{equation}
The self-adjointness of every represented element is also derived from its
signed Pauli form.

For finite $G$, define
\begin{equation}
  P_G=\frac{1}{\lvert G\rvert}\sum_{g\in G}r(g),
  \qquad
  \rho_G=2^{-n}\sum_{g\in G}r(g).
  \label{eq:finite_projector_density}
\end{equation}
The checked results give
\begin{equation}
  P_G^2=P_G=P_G^*,
  \qquad
  \Tr(P_G)=\frac{2^n}{\lvert G\rvert},
  \qquad
  \Tr(\rho_G)=1,
  \qquad
  \rho_G=\Tr(P_G)^{-1}P_G.
  \label{eq:finite_density_identities}
\end{equation}
In particular $P_G\neq0$, and $\rho_G$ is positive semidefinite.  The last
assertion is proved from the explicit self-adjoint idempotence of $P_G$, not
by an unproved positivity axiom.

\section{Finite state and code-corner trace}
\label{sec:finite_state}

The finite functional is defined by
\begin{equation}
  \omega_G(A)=\Tr(\rho_G A).
  \label{eq:finite_state_functional}
\end{equation}
Lean proves its complex linearity, normalization, and positivity on every
positive finite matrix.  In particular,
\begin{equation}
  \omega_G(\id)=1,
  \qquad
  \omega_G(P_G)=1,
  \qquad
  \omega_G(A^*A)\geq0.
  \label{eq:finite_state_properties}
\end{equation}
The functional is bundled as a positive complex-linear map for the finite
matrix order used in the proof.

The manuscript later relies on ordinary matrix-trace cyclicity after
compression to the local code space.  That step is now also checked:
\begin{equation}
  \omega_G\bigl((P_GAP_G)(P_GBP_G)\bigr)
  =
  \omega_G\bigl((P_GBP_G)(P_GAP_G)\bigr).
  \label{eq:finite_corner_traciality}
\end{equation}
The proof uses only $P_G^2=P_G$, the formula
$\rho_G=\Tr(P_G)^{-1}P_G$, and cyclicity of the finite matrix trace.  This is
the finite-dimensional precursor of the corner-traciality field in I002; it
is not yet the strong-operator extension to $q\cM q$.

\section{Signed partial trace and exact survivor condition}
\label{sec:partial_trace}

For an $n$-qubit retained subsystem and an $m$-qubit discarded subsystem, the
finite partial trace is defined entrywise.  For a signed tensor-product Pauli
term, Lean proves
\begin{equation}
  \Tr_{\mathrm{right}}\bigl(\varepsilon\,P(a)\otimes P(b)\bigr)
  =
  \begin{cases}
    2^m\,\varepsilon\,P(a),& b=(\id,\ldots,\id),\\
    0,&\text{otherwise},
  \end{cases}
  \label{eq:signed_pauli_partial_trace}
\end{equation}
where $\varepsilon\in\{+1,-1\}$.  Keeping $\varepsilon$ is essential:
consistent stabilizer groups can contain negative Hermitian Pauli strings,
even though they cannot contain $-\id$.

For any finite family of signed terms, partial trace therefore returns the
normalized signed-Pauli expression over precisely the right-identity terms.
The formal expression-compatibility structure contains the explicit equality
\begin{equation}
  T_{F'}\cap\{g:\text{$g$ is identity on }F'\setminus F\}=T_F.
  \label{eq:explicit_survivor_equality}
\end{equation}
Under \eqref{eq:explicit_survivor_equality}, the checked conclusion is
\begin{equation}
  \Tr_{F'\setminus F}(E_{F'})=E_F.
  \label{eq:checked_expression_compatibility}
\end{equation}
Here $E_{F'}$ and $E_F$ are only normalized signed-Pauli expressions.  This
is exactly the algebraic division in the manuscript calculation: Pauli trace
orthogonality proves survival, while the definition of the local stabilizer
groups must identify the survivors with $\cG_{\cS,F}$.  The current Lean
structure makes the latter requirement visible rather than silently folding it
into a generic ``stabilizer'' predicate.  It does not supply a group law,
positivity, or trace-one statement for arbitrary finite term families.

The new module \path{FiniteNestedPauli.lean} supplies the missing finite
bridge at an explicit conditional level.  Its input consists of actual finite
larger and smaller stabilizer groups, a tensor-coordinate equivalence, a
decomposition of every larger-group element into signed left/right Pauli
labels, an injective inclusion of the smaller group, and the iff-condition
that the right-identity elements are exactly the included smaller group.  It
then checks
\begin{equation}
  \Tr_{F'\setminus F}(\rho_{F'})=\rho_F,
  \label{eq:checked_actual_density_compatibility}
\end{equation}
in those tensor coordinates.  The reindexed larger density is also proved
positive semidefinite and trace-one.  Thus the generic expressions in
\eqref{eq:checked_expression_compatibility} are now identified with actual
finite density matrices under the displayed nested-group contract.  What is
not yet formalized is its full derivation from
$\cG_{\cS,F}=\cG_{\cS}\cap\cA_F$ on a countable lattice.

The checked refinement \path{FiniteLocalGroupSupport.lean} isolates one
precise part of that derivation.  Given a finite-support map with its subgroup
closure laws, together with the condition that the discarded labels are all
identity exactly when support is contained in the smaller region, it defines
the two local groups by support containment.  It then derives their canonical
inclusion and the survivor iff-condition used above.  Thus those two clauses
are no longer independent fields once such support data is supplied.  The
support realization itself, the Pauli-coordinate factorization, and the
discarded-label detection condition remain inputs to be constructed from the
actual countable-lattice Pauli group and its local algebras.

The algebraic normal-form modules now supply this discarded-label detection
for their own syntactic support, when the finite complement is explicitly
enumerated.  More precisely, the projective $X/Z$ coordinates are assigned
the canonical Hermitian labels $\id,X,Y,Z$, and a site belongs to the support
exactly when its label is not $\id$.  This records the survivor calculation
for an enumerated complement without an extra hypothesis.  It does not
represent the central $\mathbb Z/4$ phase of a full normal-form word as a
matrix phase, so it is not a realization of the full Pauli group or of the
manuscript's local matrix algebra.

The finite-coordinate refinement makes the restriction calculation usable in
the physical ordered-qubit convention.  The checked equivalence between the
recursive computational basis and functions on $\operatorname{Fin}(n)$ splits
the first $n$ and final $m$ labels explicitly.  The tensor inclusion is then
the reindexed map $A\mapsto A\otimes\id$, and Lean checks its trace adjointness
with the right partial trace.  Combining this with
\eqref{eq:checked_actual_density_compatibility} proves equality of the actual
finite state functionals on one inclusion.  Longer maps are defined by
recursion from those physical one-site embeddings, and a stepwise compatible
family has the checked equality of state values along every displayed
sequential transition.

The module \path{SequentialLocalFunctional.lean} then uses those equalities to
define a complex-linear functional on the corresponding
\texttt{Module.DirectLimit}.  The finite component maps are checked injective,
operator-norm-preserving unital star-algebra homomorphisms with identity and
composition laws, so the quotient is now a module direct limit of a genuine
directed system.  The
endpoint-to-added-block bridge transfers the actual signed-Pauli padding
formula to these typed maps.  Mathlib's module-colimit construction still
supplies no multiplication or star operation on this noncommutative quotient.
Positivity there and norm completion remain separate obligations.  Accordingly
the new quotient is not called a quasi-local $C^*$-algebra and its functional
is not called a global state.

\section{Bridge to the support-corner theorem}
\label{sec:bridge}

The finite stabilizer representation can be transported through any unital
star-algebra representation.  The checked image of the group average is the
image of $P_G$.  Consequently, once an actual quasi-local GNS representation
has supplied the represented local stabilizers, this finite construction is
the same group-average input consumed by the decreasing-support theorem of
I002.  It is not a second, incompatible notion of local support.

The separate package module \path{GNSVector.lean} also supplements
mathlib's generic GNS construction.  For any supplied positive linear
functional on a unital $C^*$-algebra, the image of the unit reconstructs that
functional as a vector coefficient and has dense represented-algebra orbit;
it has norm one if the functional is normalized.  This does not construct the
compatible quasi-local functional from the finite densities or prove that the
represented local stabilizers fix its vector.

At finite volume, \path{FinitePauliGNS.lean} applies this construction to the
actual stabilizer density.  It derives that each represented stabilizer has
zero squared deviation in its own finite state and hence fixes the resulting
GNS vector.  This is a finite-dimensional statement; it does not provide the
compatible countable-lattice state needed to identify a global GNS vector.

The following dependencies are now checked:
\begin{equation}
  \begin{aligned}
    \text{signed Pauli data}
    &\longrightarrow (P_G,\rho_G,\omega_G) \\
    &\longrightarrow \text{finite code-corner traciality} \\
    &\longrightarrow \text{represented finite group averages of I002}.
  \end{aligned}
  \label{eq:checked_dependency_chain}
\end{equation}

\section{Remaining boundary}
\label{sec:boundary}

The following are not claimed by I003.
\begin{enumerate}[leftmargin=2.2em]
  \item Construction of the global Pauli group on a countable lattice, its
  concrete local algebras, and a finite-support realization identifying
  support containment with the global definition
  $\cG_{\cS,F}=\cG_{\cS}\cap\cA_F$ for arbitrary finite lattice regions.
  \item A full matrix realization of the $\mathbb Z/4$-phased normal form and
  its relation to arbitrary finite lattice regions.  The literal ordered-block
  factorization and syntactic discarded-label test for an enumerated complement
  are checked, but they do not identify support with membership in a global
  local matrix algebra.
  \item Construction of a noncommutative star-algebra colimit on the checked
  coherent finite system, followed by a norm-closed quasi-local $C^*$-algebra
  and extension of \eqref{eq:finite_state_functional} to a positive global
  state.  The present module direct limit and its functional are deliberately
  weaker than this.
  \item Construction of the compatible quasi-local functional whose generic
  GNS vector is supplied by \path{GNSVector.lean}, and a compatibility proof
  lifting the finite \path{FinitePauliGNS.lean} fixed-vector calculation to
  represented local $P_G$.
  \item Extension of \eqref{eq:finite_corner_traciality} from finite local
  operators to all of $q\cM q$ through bounded strong convergence and
  Kaplansky density.
  \item Proof of the regional/complementary product-orbit density used for
  fullness and projection-faithfulness in I002.
\end{enumerate}

Thus I003 formalizes the actual finite-group density input, the signed
partial-trace algebra with its correct signs and survivor hypothesis, the
conditional equality of nested local stabilizer density matrices, and
finite-volume sequential state compatibility.  It also supplies a linear
functional on the module direct limit of its coherent finite star-algebra
system, with intentionally limited scope.  The infinite-volume conclusion
remains the conditional, explicitly parameterized theorem of I002 until the
six listed steps are implemented.

\section{Verification record}
\label{sec:verification}

The intended Lean verification is
\begin{equation}
  \texttt{lake build}
  \qquad\text{and}\qquad
  \texttt{lake env lean tests/AxiomAudit.lean}
  \label{eq:lean_verification_commands}
\end{equation}
from \texttt{formalization\_finite\_fibers/}.  The latter audit prints only
Lean's standard logical axioms for the public theorems and no project-local
axiom.  Independent mathematical-fidelity and Lean/API adversarial reviews
for this investigation are recorded alongside this note.

\end{document}
